% ============ 第七章 §3 离散赋值域 (原书 pp.254-260) ============
\section{离散赋值域}\label{sec:ch7-s3}

本节对非阿赋值域, 特别是对一类特殊的非阿赋值域——离散赋值域作进一步研究. 超距不等式 $\varphi(a+b)\leqslant\max(\varphi(a),\varphi(b))$ 使得非阿赋值域有许多特别的性质.

设 $(F,\varphi)$ 是非阿赋值域, 对每个 $a\in F^{\times}$, 令
\[
  v(a)=-\log\varphi(a),\qquad\text{即 } \varphi(a)=e^{-v(a)},
\]
而令 $v(0)=\infty$, 则 $v$ 是从 $F$ 到集合 $\R\cup\{\infty\}$ 的映射, 并且非阿赋值 $\varphi$ 的 $3$ 条性质可转化成 $v$ 的如下 $3$ 条性质: 对 $a,b\in F$,
\[
  (1)\ v(a)=\infty\ \text{当且仅当 } a=0,
\]
\[
  (2)\ v(ab)=v(a)+v(b),
\]
\[
  (3)\ v(a+b)\geqslant\min(v(a),v(b)).
\]
我们把满足这 $3$ 条性质的 $v$ 叫作域 $F$ 的\emph{指数赋值}. 反过来, 给了 $F$ 的指数赋值 $v$, 则对每个实数 $\gamma\ (0<\gamma<1)$, $\varphi(a)=\gamma^{v(a)}$ 是 $F$ 的非阿赋值. 所以 $F$ 的非阿赋值和指数赋值相互对应, 从而一个非阿赋值域也可写成 $(F,v)$. 进而, 若 $\varphi_1(a)=\gamma^{v_1(a)}$, $\varphi_2(a)=\gamma^{v_2(a)}$, 则 $\varphi_1=\varphi_2^{\alpha}\Leftrightarrow v_1=\alpha v_2$. 所以若把域 $F$ 的两个指数赋值 $v_1$ 和 $v_2$ 等价定义为存在正实数 $\alpha$, 使得 $v_1=\alpha v_2$, 那么等价的赋值对应等价的指数赋值. 所以域 $F$ 的一个非阿素除子也可说成是一个指数赋值等价类.

设 $P$ 是域 $F$ 的非阿素除子, $v$ 是 $P$ 中的一个指数赋值. 由指数赋值的性质 (2) 可知
\[
  v\colon F^{\times}\to\R
\]
是乘法群 $F^{\times}$ 到实数加法群 $\R$ 的同态, 所以集合 $v(F^{\times})$ 为 $\R$ 的加法子群.

\begin{definition}\label{def:ch7-5}
非阿指数赋值 $v$ 叫作\emph{离散的}, 是指 $v(F^{\times})$ 是 $\R$ 的离散子群 $\gamma\Z$, 其中 $\gamma$ 是某个正实数.
\end{definition}

若 $v$ 是离散的, 则与之等价的每个指数赋值 $\alpha v\ (\alpha>0)$ 也是离散的（因为 $\alpha v(F^{\times})=\alpha\gamma\Z$）. 所以把 $v$ 所在的非阿素除子 $P$ 也叫作离散的, 并且 $P$ 中的非阿赋值 $\varphi$ 也叫作离散的, 而域 $F$（或者更确切地说应当是 $(F,P)$, $(F,\varphi)$ 或者 $(F,v)$）叫作\emph{离散赋值域}. 注意: 只有对非阿素除子和非阿（指数）赋值才有离散素除子和离散（指数）赋值概念, 它与``平凡赋值给出离散拓扑''这一命题中的``离散''是两件不同的事情.

设 $v$ 为域 $F$ 的离散指数赋值, 则 $v(F^{\times})=\gamma\Z\ (\gamma>0)$. 取 $v'=\gamma^{-1}v$ 为与 $v$ 等价的指数赋值, 则 $v'(F^{\times})=\Z$. 我们把 $v'$ 叫作离散素除子 $P$ 中的\emph{标准指数赋值}, 通常表示成 $v_P$. 比如说: 第 $7.1$ 节中在有理数域 $\Q$ 的 $p$ 进指数赋值 $v_p$ 就是标准的离散指数赋值, 因为 $v_p(\Q^{\times})=\Z$.

下面是非阿赋值域的一个重要性质:

\begin{lemma}\label{lem:ch7-6}
设 $(F,P)$ 是非阿赋值域, $\varphi$ 和 $v$ 分别是素除子 $P$ 中的赋值和指数赋值, 令
\[
  O=\{a\in F\mid\varphi(a)\leqslant 1\}=\{a\in F\mid v(a)\geqslant 0\},
\]
\[
  \mathfrak{p}=\{a\in F\mid\varphi(a)<1\}=\{a\in F\mid v(a)>0\},
\]
\[
  U=O-\mathfrak{p}=\{a\in F\mid\varphi(a)=1\}=\{a\in F\mid v(a)=0\},
\]
则 $O$ 是 $F$ 的子环, $\mathfrak{p}$ 为环 $O$ 的素理想, $U$ 是环 $O$ 的单位群, 并且 $F$ 是 $O$ 的商域（即 $F$ 中元素均可表为 $O$ 中两个元素的商）.
\end{lemma}

\begin{proof}
由超距不等式可知 $O$ 是环, 其余可由定义直接推出.
\end{proof}

\begin{remark}
(1) 设 $\varphi_1$ 和 $\varphi_2$ 是素除子 $P$ 中两个赋值, 则 $\varphi_1(a)\leqslant 1\Leftrightarrow\varphi_2(a)\leqslant 1$. 因此由 $\varphi_1$ 和 $\varphi_2$ 定义出同一个环 $O$, 即 $O$ 只是非阿素除子 $P$ 的特性. 同样地, $\mathfrak{p}$ 和 $U$ 也是素除子 $P$ 的特性. 我们把 $O$, $\mathfrak{p}$ 和 $U$ 分别叫作 $F$ 对于非阿素除子 $P$ 的\emph{赋值环}, \emph{素理想}和\emph{单位群}, 而 $O$ 有时也表示成 $O_P$.

(2) 由于 $O=\mathfrak{p}\cup U$, $\mathfrak{p}\cap U=\varnothing$, 可知 $\mathfrak{p}$ 事实上为环 $O$ 的惟一极大理想, 所以 $O$ 是局部环. 特别地, $O/\mathfrak{p}$ 是域, 叫作 $F$ 在 $P$ 的\emph{剩余类域}, 表示成 $\overline{F}$.
\end{remark}

\begin{example}
设 $p$ 为素数, 则非阿赋值域 $(\Q,v_p)$ 的赋值环, 素理想和单位群分别为
\begin{align*}
  O&=O_p=\{\alpha\in\Q\mid v_p(\alpha)\geqslant 0\}\\
  &=\left\{\alpha=\frac{m}{n}\mid m,n\in\Z,\ (p,n)=1\right\},\\
  \mathfrak{p}&=\{\alpha\in\Q\mid v_p(\alpha)>0\}\\
  &=\left\{\alpha=\frac{m}{n}\mid m,n\in\Z,\ (p,n)=1,\ p\mid m\right\}=pO,\\
  U&=\{\alpha\in\Q\mid v_p(\alpha)=0\}\\
  &=\left\{\alpha=\frac{m}{n}\mid m,n\in\Z,\ (mn,p)=1\right\}=O-pO,
\end{align*}
而剩余类域为 $O/\mathfrak{p}\cong\Z/p\Z$（$p$ 元域）.
\end{example}

\emph{离散赋值域的结构} 设 $(F,P)$ 为离散赋值域, $\varphi$ 和 $v$ 分别为 $P$ 中的非阿赋值和标准指数赋值. 于是 $v(F^{\times})=\Z$, 从而存在元素 $\pi\in F^{\times}$, 使得 $v(\pi)=1$. 这样的元素 $\pi$ 叫作 $F$（对于 $P$）的\emph{素元}. $F$ 中不同的素元彼此相差一个单位因子（$\pi'$ 为素元 $\Leftrightarrow v(\pi')=1\Leftrightarrow v\left(\frac{\pi'}{\pi}\right)=0\Leftrightarrow\frac{\pi'}{\pi}\in U$）. 不难看出 $\pi\in\mathfrak{p}$, 并且 $\mathfrak{p}=\pi O$. 所以 $\mathfrak{p}$ 是环 $O$ 中的主理想. 更一般的, 设 $A$ 为 $O$ 中的非零理想, 令
\[
  r=\min\{v(a)\mid 0\neq a\in A\}.
\]
由 $A\neq(0)$ 可知 $r\geqslant 0$, 并且存在 $a\in A$, 使得 $v(a)=r$, 于是 $\pi^r/a\in U$, 从而 $\pi^r\in aU\subseteq A$, 即 $\pi^rO\subseteq A$. 反之, 对每个 $b\in A$, $v(b)\geqslant r$, 所以 $v(b/\pi^r)\geqslant 0$, 即 $b/\pi^r\in O$, 从而 $b\in\pi^rO$. 于是 $A\subseteq\pi^rO$. 这就证明了 $O$ 的每个非零理想均是主理想 $A=\pi^rO=\mathfrak{p}^r\ (r=0,1,2,\cdots)$. 这些理想（对不同的 $r$）彼此不同, 并且全体构成零元素 $0$ 的一组基本开集系. 由于
\[
  \mathfrak{p}^r=\{a\in F\mid v(a)\geqslant r\}=\{a\in F\mid v(a)>r-1\},
\]
可知对于拓扑 $T(P)$, 每个理想 $\mathfrak{p}^r$ 都是又开又闭的, 所以离散赋值域的拓扑是全不连通的, 即 $F$ 的每个点都是一个连通分支.

$F$ 中每个非零元素 $a$ 均可表成
\[
  a=\pi^ru,\qquad r=v(a),\ u\in U,
\]
所以乘法群 $F^{\times}$ 是由 $\pi$ 生成的自由循环群 $\langle\pi\rangle$ 和单位群 $U$ 的直积.

$F$ 的分式理想是指 $F$ 的一个子集 $M$, 并且存在 $a\in F^{\times}$, 使得 $aM$ 为 $O$ 的非零理想. 由上述知 $aM=\pi^nO\ (n\geqslant 0)$, 令 $a=\pi^ru\ (r\in\Z,\ u\in U)$, 则
\[
  M=a^{-1}\pi^nO=u^{-1}\pi^{n-r}O=\pi^{n-r}O=\mathfrak{p}^{n-r},
\]
其中 $n-r\in\Z$. 反之, 对每个 $m\in\Z$, $\mathfrak{p}^m=\pi^mO$ 都是 $F$ 的分式理想（因为 $\pi^{-m}\mathfrak{p}^m=O$）. 这就表明: $F$ 的全部分式理想为 $\mathfrak{p}^n\ (n\in\Z)$, 并且 $\mathfrak{p}^n$ 为 $O$ 的理想当且仅当 $n\geqslant 0$.

\emph{离散赋值域 $F$ 的加法结构} 我们上面给出 $F$ 的一个加法子群链
\[
  F\supset\cdots\supset\mathfrak{p}^{-2}\supset\mathfrak{p}^{-1}\supset\mathfrak{p}^0=O\supset\mathfrak{p}\supset\mathfrak{p}^2\supset\cdots\supset(0),
\]
不难看出, 对任意两个整数 $r$ 和 $s$, 映射
\[
  \mathfrak{p}^r\to\mathfrak{p}^s,\qquad a\mapsto\pi^{s-r}a
\]
是加法群同构（事实上这是拓扑群同构, 即上述映射是连续的）. 另一方面, 设 $r\in\Z$ 而 $n$ 为非负整数, 则映射
\[
  f\colon\mathfrak{p}^r\to O/\mathfrak{p}^n,\qquad a\mapsto a\pi^{-r}\ (\mathrm{mod}\ \mathfrak{p}^n)
\]
为加法群满同态, 并且
\begin{align*}
  a\in\operatorname{Ker}(f)&\Leftrightarrow a\pi^{-r}\in\mathfrak{p}^n\Leftrightarrow v(a\pi^{-r})\geqslant n\\
  &\Leftrightarrow v(a)\geqslant r+n\Leftrightarrow a\in\mathfrak{p}^{n+r},
\end{align*}
于是 $\operatorname{Ker}(f)=\mathfrak{p}^{n+r}$, 所以有加法群同构
\[
  \mathfrak{p}^r/\mathfrak{p}^{r+n}\cong O/\mathfrak{p}^n\qquad(r,n\in\Z,\ n\geqslant 0).
\]
特别地, $\mathfrak{p}^r/\mathfrak{p}^{r+1}\cong O/\mathfrak{p}=\overline{F}$. 如果 $\overline{F}$ 是有限域 $\F_q$, 则
\begin{align*}
  [\mathfrak{p}^r:\mathfrak{p}^{r+n}]&=[\mathfrak{p}^r:\mathfrak{p}^{r+1}][\mathfrak{p}^{r+1}:\mathfrak{p}^{r+2}]\cdots[\mathfrak{p}^{r+n-1}:\mathfrak{p}^{r+n}]\\
  &=[O:\mathfrak{p}]^n=|\overline{F}|^n=q^n.
\end{align*}
特别地, 对每个 $n\geqslant 1$, $|O/\mathfrak{p}^n|=q^n$.

对于 $a,b\in F$, $r\in\Z$, 我们用同余式
\[
  a\equiv b\ (\mathrm{mod}\ \mathfrak{p}^r)
\]
表示 $a-b\in\mathfrak{p}^r$（即 $v(a-b)\geqslant r$）, 当 $r\geqslant 1$ 时, 这与通常环 $O$ 中理想 $\mathfrak{p}^r$ 的同余符号的含义是一致的.

\emph{$F^{\times}$ 的乘法结构} 乘法拓扑群 $F^{\times}$ 有如下一些子群: $U_0=U$, 而对 $r\geqslant 1$,
\[
  U_r=1+\mathfrak{p}^r=\{a\in F^{\times}\mid v(a-1)\geqslant r\}
  =\{a\in F^{\times}\mid v(a-1)>r-1\}
\]
（请验证 $U_r$ 是子群）, 于是我们有 $F^{\times}$ 的乘法子群链
\[
  F^{\times}\supset U_0\supset U_1\supset\cdots\supset\{1\},
\]
它们形成 $1$ 的又开又闭的基本开集系, 所以乘法群 $F^{\times}$ 也是全不连通拓扑群.

\begin{lemma}\label{lem:ch7-7}
(1) $F^{\times}/U_0\cong\Z$（整数加法群）;

(2) $U_0/U_1\cong\overline{F}^{\times}$（乘法群）;

(3) 当 $r\geqslant 1$ 时, $U_r/U_{r+1}\cong\overline{F}$（加法群）.
\end{lemma}

\begin{proof}
(1) 考虑标准指数赋值映射
\[
  v\colon F^{\times}\to\Z,\qquad a\mapsto v(a),
\]
由 $v(\pi)=1$ 可知这是群的满同态. $\operatorname{Ker}(v)=U_0$, 于是 $F^{\times}/U_0\cong\Z$.

(2) 注意 $\overline{F}=O/\mathfrak{p}$, $U=O-\mathfrak{p}$. 考虑映射
\[
  f\colon U_0\to\overline{F}^{\times},\qquad a\mapsto\bar{a}\quad(\bar{a}\ \text{表示 } a\ (\mathrm{mod}\ \mathfrak{p})),
\]
这是乘法满同态, 并且
\[
  a\in\operatorname{Ker}(f)\Leftrightarrow\bar{a}=\overline{1}\Leftrightarrow a\equiv 1\ (\mathrm{mod}\ \mathfrak{p})\Leftrightarrow a\in U_1,
\]
从而 $U_0/U_1\cong\overline{F}^{\times}$.

(3) 考虑映射
\[
  f\colon U_r=1+\mathfrak{p}^r\to\overline{F},\qquad 1+\pi^ra\mapsto\bar{a}\quad(a\in O).
\]
由于 $r\geqslant 1$, 可知对 $a,b\in O$,
\begin{align*}
  f((1+\pi^ra)(1+\pi^rb))&=f(1+\pi^r(a+b+\pi^rab))\\
  &=\overline{a+b+\pi^rab}=\overline{a+b}=f(1+\pi^ra)+f(1+\pi^rb),
\end{align*}
从而 $f$ 是群同态, 并且显然是满同态. 进而对 $a\in O$,
\[
  1+\pi^ra\in\operatorname{Ker}(f)\Leftrightarrow\bar{a}=\overline{0}\Leftrightarrow a\in\mathfrak{p}
  \Leftrightarrow 1+\pi^ra\in 1+\mathfrak{p}^{r+1}=U_{r+1},
\]
从而 $U_r/U_{r+1}\cong\overline{F}$.
\end{proof}

\begin{remark}
如果 $\overline{F}$ 为有限域 $\F_q$, 则当 $r\geqslant 1$ 时,
\begin{align*}
  |U_0/U_r|&=|U_0/U_1|\cdot|U_1/U_2|\cdots|U_{r-1}/U_r|\\
  &=(q-1)q^{r-1}.
\end{align*}

对于 $r\geqslant 1$, $a,b\in F^{\times}$, 我们用符号
\[
  a\equiv b\ (\mathrm{mod}^{\times}\mathfrak{p}^r)
\]
表示 $a/b\in U_r$（即 $v\left(\frac{a}{b}-1\right)\geqslant r$）. 并且规定
\[
  a\equiv b\ (\mathrm{mod}^{\times}\mathfrak{p}^0)
\]
表示 $a/b\in U_0=U$（即 $v(a)=v(b)$）.
\end{remark}

\begin{Exercises}
\item 设 $F$ 为域, $F(x)$ 是域 $F$ 上关于 $x$ 的有理函数域, $F[x]$ 是域上关于 $x$ 的多项式环.

(1) 设 $p=p(x)$ 为 $F[x]$ 中最高项系数为 $1$ 的 $d$ 次不可约多项式（$d\geqslant 1$）, 对 $F[x]$ 中每个非零多项式 $a=a(x)$, 令 $p^n\parallel a$（即 $p^n\mid a$, $p^{n+1}\nmid a$）. 我们定义 $v_p(a)=n$, 而对每个非零有理函数 $\alpha=a(x)/b(x)$（$a(x),b(x)\in F[x]$, $b(x)\neq 0$）, 定义 $v_p(\alpha)=v_p(a(x))-v_p(b(x))$. 再令 $v_p(0)=\infty$. 求证 $v_p$ 是域 $F(x)$ 的标准离散指数赋值.

(2) 描述离散赋值域 $(F(x),v_p)$ 的赋值环 $O$, 素理想 $\mathfrak{p}$ 和单位群 $U$. 证明 $\overline{F(x)}=O/\mathfrak{p}\cong F(\theta)$, 其中 $\theta$ 是 $p(x)$ 在 $F$ 的代数闭包中的一个根（从而 $[F(\theta):F]=d$）, $v_p$ 所在的 $F(x)$ 的素除子也简记为 $p$.

\item $F$, $F(x)$, $F[x]$ 如前题所示, 对每个非零有理函数 $\alpha=a(x)/b(x)$（$a(x),b(x)\in F[x]$, $b(x)\neq 0$）, 定义
\[
  v_{\infty}(\alpha)=\deg b(x)-\deg a(x),
\]
其中 $\deg f(x)$ 表示多项式 $f(x)$ 的次数, 并且规定 $v_{\infty}(0)=\infty$, $\deg(0)=-\infty$.

(1) 求证 $v_{\infty}$ 是域 $F(x)$ 的标准离散指数赋值, $\frac{1}{x}$ 为对 $v_{\infty}$ 的一个素元, $v_{\infty}$ 所在的素除子也记为 $\infty$.

(2) 描述离散赋值域 $(F(x),v_{\infty})$ 的赋值环 $O$, 素理想 $\mathfrak{p}$ 和单位群 $U$. 证明 $O/\mathfrak{p}\cong F$.

(3) 证明对 $F[x]$ 中不同的首 $1$ 不可约多项式 $p_1$ 和 $p_2$, $v_{p_1}$ 和 $v_{p_2}$ 不等价, 并且也和 $v_{\infty}$ 不等价.

\item 设 $v$ 为域 $F$ 的指数赋值, 其赋值环和素理想分别为 $O$ 和 $\mathfrak{p}$, $\overline{F}=O/\mathfrak{p}$.

(1) 对于多项式 $f(x)=a_0+a_1x+\cdots+a_nx^n\in F[x]$, 定义
\[
  \bar{v}(f)=\min_{0\leqslant i\leqslant n}\{v(a_i)\}.
\]
对于 $\alpha=f/g\in F(x)$（$f,g\in F[x]$, $g\neq 0$）, 定义 $\bar{v}(\alpha)=\bar{v}(f)-\bar{v}(g)$. 求证 $\bar{v}$ 是域 $F(x)$ 的指数赋值, 并且 $\bar{v}$ 在 $F$ 上的限制就是 $v$.

(2) 设 $f(x)$ 是 $O[x]$ 中的首 $1$ 多项式, $f(x)=g(x)h(x)$, 其中 $g(x)$ 和 $h(x)$ 均是 $F[x]$ 中的首 $1$ 多项式, 求证 $g(x),h(x)\in O[x]$.

(3) 多项式 $f(x)\in O[x]$ 叫作\emph{本原的}, 是指 $\bar{v}(f)=0$. 求证: 两个本原多项式的乘积仍是本原多项式.

(4) 试刻画 $F(x)$ 对于 $\bar{v}$ 的赋值环 $O'$ 和素理想 $\mathfrak{p}'$. 证明 $O'/\mathfrak{p}'\cong\overline{F}(x)$.

(5) 设 $f(x)$ 是 $O[x]$ 中首 $1$ 多项式, 则: $f(x)$ 在 $O[x]$ 中不可约 $\Leftrightarrow f(x)$ 在 $F[x]$ 中不可约.

(6) （Eisenstein 判别法）设 $f(x)=x^n+a_{n-1}x^{n-1}+\cdots+a_0\in O[x]$, $a_0,\cdots,a_{n-1}\in\mathfrak{p}$, 并且 $a_0$ 不是 $\mathfrak{p}$ 中两个元素之积. 求证 $f(x)$ 在 $F[x]$ 中不可约.

\item 设 $(F,v)$ 是离散赋值域, $v$ 为标准指数赋值, $\pi$ 为一个素元, $r\geqslant 1$. 如果
\[
  a=\pi^{v(a)}\varepsilon,\qquad b=\pi^{v(b)}\eta,\qquad\varepsilon,\eta\in U,
\]
则 $a\equiv b\ (\mathrm{mod}^{\times}\mathfrak{p}^r)\Leftrightarrow v(a)=v(b)$ 并且 $\varepsilon\equiv\eta\ (\mathrm{mod}\ \mathfrak{p}^r)$.

\item 设 $(F,P)$ 是离散赋值域, $O$ 和 $\mathfrak{p}$ 为其赋值环和素理想, 域 $\overline{F}=O/\mathfrak{p}$ 的特征为素数 $p$, $r\geqslant 1$.

(1) 若 $a,b\in O$, $a\equiv b\ (\mathrm{mod}\ \mathfrak{p}^r)$, 则对每个 $s\geqslant 0$,
\[
  a^{p^s}\equiv b^{p^s}\ (\mathrm{mod}\ \mathfrak{p}^{r+s}).
\]
(2) 若 $a,b\in F^{\times}$, $a\equiv b\ (\mathrm{mod}^{\times}\mathfrak{p}^r)$, 则对每个 $s\geqslant 0$,
\[
  a^{p^s}\equiv b^{p^s}\ (\mathrm{mod}^{\times}\mathfrak{p}^{r+s}).
\]
\end{Exercises}
